% =====================================================================
% CIRCUITOS REVISADOS --- ANALISE NODAL E ANALISE DE MALHAS
% Correcoes visuais de circuitos com sobreposicao/espacamento insuficiente.
% =====================================================================

% ----------------------------- NODAL ----------------------------------
\newcommand{\fvNodalPonteRevisada}{%
\begin{circuitikz}[american,scale=.88,transform shape,line width=.8pt]
  \coordinate (g) at (0,0);
  \draw (g) to[isource,l_={$\SI{6}{\ampere}$}] (0,3)
        -- (1.4,3) coordinate(n1);
  \draw (n1) to[R,l={$\SI{4}{\ohm}$}] (4.6,4.2) coordinate(n2);
  \draw (n1) to[R,l_={$\SI{6}{\ohm}$}] (4.6,1.8) coordinate(n3);
  \draw (n2) to[R,l={$\SI{10}{\ohm}$}] (n3);
  \draw (n2) to[R,l={$\SI{8}{\ohm}$}] (7.0,0) coordinate(gr);
  \draw (n3) to[R,l_={$\SI{12}{\ohm}$}] (4.6,0);
  \draw (gr) -- (g);
  \node[ground] at(g){};
  \node[Vcol,fill=white,inner sep=1pt] at ($(n1)+(0,.35)$) {$V_1$};
  \node[Vcol,fill=white,inner sep=1pt] at ($(n2)+(0,.35)$) {$V_2$};
  \node[Vcol,fill=white,inner sep=1pt] at ($(n3)+(.35,0)$) {$V_3$};
\end{circuitikz}}

\newcommand{\fvNodalPortaTheveninRevisada}{%
\begin{circuitikz}[american,scale=.9,transform shape,line width=.8pt]
  \draw (0,0) to[isource,l_={$\SI{4}{\ampere}$}] (0,3)
        -- (1.5,3) coordinate(n1);
  \draw (n1) to[R,l={$\SI{10}{\ohm}$}] (1.5,0);
  \draw (n1) to[R,l={$\SI{20}{\ohm}$}] (4.9,3)
        node[ocirc]{} node[above=3pt,Vcol] {$a$};
  \draw (1.5,0)--(0,0);
  \node[ground] at(1.5,0){};
  \node[Vcol,fill=white,inner sep=1pt] at ($(n1)+(0,.35)$) {$V_1$};
\end{circuitikz}}

\newcommand{\fvNodalEscadaQuatroRevisada}{%
\begin{circuitikz}[american,scale=.84,transform shape,line width=.8pt]
  \draw (0,0) to[isource,l_={$\SI{4}{\ampere}$}] (0,3)
        -- (1.25,3) coordinate(n1);
  \draw (n1) to[R,l_={$\SI{10}{\ohm}$}] (3.55,3) coordinate(n2);
  \draw (n2) to[R,l_={$\SI{10}{\ohm}$}] (5.85,3) coordinate(n3);
  \draw (n3) to[R,l_={$\SI{10}{\ohm}$}] (8.15,3) coordinate(n4);
  \draw (n1) to[R,l={$\SI{10}{\ohm}$}] (1.25,0);
  \draw (n2) to[R,l={$\SI{10}{\ohm}$}] (3.55,0);
  \draw (n3) to[R,l={$\SI{10}{\ohm}$}] (5.85,0);
  \draw (n4) to[R,l={$\SI{10}{\ohm}$}] (8.15,0) -- (0,0);
  \node[ground] at(1.25,0){};
  \foreach \p/\lab in {n1/1,n2/2,n3/3,n4/4}{%
    \node[Vcol,fill=white,inner sep=1pt] at ($(\p)+(0,.38)$) {$V_\lab$};}
\end{circuitikz}}

\newcommand{\fvNodalPortaDependenteRevisada}{%
\begin{circuitikz}[american,scale=.88,transform shape,line width=.8pt]
  \draw (0,0) to[isource,l_={$\SI{2}{\ampere}$}] (0,3)
        -- (1.5,3) coordinate(n1);
  \draw (n1) to[R,l={$\SI{10}{\ohm}$}] (1.5,0);
  \draw (n1) to[R,l={$\SI{5}{\ohm}$}] (4.7,3) coordinate(n2);
  \draw (n2) -- (6.0,3) node[ocirc]{} node[above=3pt,Vcol] {$a$};
  \draw (7.8,0) to[cisource,l_={$0{,}05V_1$}] (7.8,3) -- (n2);
  \draw (7.8,0)--(0,0);
  \node[ground] at(1.5,0){};
  \node[Vcol,fill=white,inner sep=1pt] at ($(n1)+(0,.35)$) {$V_1$};
  \node[Vcol,fill=white,inner sep=1pt] at ($(n2)+(0,.35)$) {$V_2$};
\end{circuitikz}}

% ----------------------------- MALHAS ---------------------------------
\newcommand{\fvMalhaSuperRevisada}{%
\begin{circuitikz}[american,scale=.95,transform shape,line width=.8pt]
  \draw (0,0) to[battery1,l^={$\SI{24}{\volt}$}] (0,3.2)
        to[R,l={$R_1=\SI{4}{\ohm}$}] (3.2,3.2) coordinate(t);
  \draw (t) to[isource,l={$\SI{2}{\ampere}$}] (3.2,0) coordinate(b) -- (0,0);
  \draw (t) to[R,l={$R_2=\SI{4}{\ohm}$}] (6.4,3.2) -- (6.4,0) -- (b);
  \draw[->,loopcol,line width=1.2pt] (1.55,1.25) arc (0:305:.38);
  \node[loopcol,fill=white,inner sep=1pt] at(1.55,1.25){\footnotesize$I_1$};
  \draw[->,loopcol,line width=1.2pt] (4.85,1.25) arc (0:305:.38);
  \node[loopcol,fill=white,inner sep=1pt] at(4.85,1.25){\footnotesize$I_2$};
  \draw[dashed,laranja,line width=.9pt,rounded corners=7pt]
        (-.75,-.6) rectangle (7.15,3.85);
  \node[laranja,font=\footnotesize\sffamily,fill=white,inner sep=2pt]
        at (3.2,-.95) {supermalha};
\end{circuitikz}}

\newcommand{\fvMalhaTresRevisada}{%
\begin{circuitikz}[american,scale=.9,transform shape,line width=.8pt]
  \draw (0,0) to[battery1,l^={$\SI{24}{\volt}$}] (0,3.2)
        to[R,l={$R_1=\SI{4}{\ohm}$}] (3.1,3.2) coordinate(t1);
  \draw (t1) to[R,l_={$R_2=\SI{4}{\ohm}$}] (3.1,0) coordinate(b1)--(0,0);
  \draw (t1) to[R,l={$R_3=\SI{2}{\ohm}$}] (6.2,3.2) coordinate(t2);
  \draw (t2) to[R,l_={$R_4=\SI{4}{\ohm}$}] (6.2,0) coordinate(b2)--(b1);
  \draw (t2) to[R,l={$R_5=\SI{4}{\ohm}$}] (9.3,3.2)--(9.3,0)--(b2);
  \foreach \x/\n in {1.45/1,4.65/2,7.85/3}{%
    \draw[->,loopcol,line width=1.1pt] (\x,1.15) arc(0:305:.36);
    \node[loopcol,fill=white,inner sep=1pt] at(\x,1.15){\footnotesize$I_\n$};}
\end{circuitikz}}

\newcommand{\fvMalhaTresDuasFontesRevisada}{%
\begin{circuitikz}[american,scale=.9,transform shape,line width=.8pt]
  \draw (0,0) to[battery1,l^={$E_1=\SI{24}{\volt}$}] (0,3.2)
        to[R,l={$\SI{2}{\ohm}$}] (3.1,3.2) coordinate(t1);
  \draw (t1) to[R,l_={$\SI{2}{\ohm}$}] (3.1,0) coordinate(b1)--(0,0);
  \draw (t1) to[R,l={$\SI{2}{\ohm}$}] (6.2,3.2) coordinate(t2);
  \draw (t2) to[R,l_={$\SI{2}{\ohm}$}] (6.2,0) coordinate(b2)--(b1);
  \draw (t2) to[R,l={$\SI{6}{\ohm}$}] (9.3,3.2)
        to[battery1,l_={$E_2=\SI{12}{\volt}$},invert] (9.3,0)--(b2);
  \foreach \x/\n in {1.45/1,4.65/2,7.85/3}{%
    \draw[->,loopcol,line width=1.1pt] (\x,1.15) arc(0:305:.36);
    \node[loopcol,fill=white,inner sep=1pt] at(\x,1.15){\footnotesize$I_\n$};}
\end{circuitikz}}

\newcommand{\fvMalhaQuatroRevisada}{%
\begin{circuitikz}[american,scale=.9,transform shape,line width=.8pt]
  \draw (0,2.8) to[battery1,l^={$E=\SI{48}{\volt}$}] (0,5.6)
        to[R,l={$\SI{4}{\ohm}$}] (3.5,5.6) coordinate(A);
  \draw (A) to[R,l={$\SI{2}{\ohm}$}] (7.0,5.6) coordinate(B);
  \draw (A) to[R,l_={$\SI{2}{\ohm}$}] (3.5,2.8) coordinate(C);
  \draw (B) to[R,l={$\SI{2}{\ohm}$}] (7.0,2.8) coordinate(D);
  \draw (0,2.8)--(C)--(D);
  \draw (0,2.8) to[R,l^={$\SI{1}{\ohm}$}] (0,0) coordinate(E);
  \draw (E) to[R,l_={$\SI{2}{\ohm}$}] (3.5,0) coordinate(F);
  \draw (F) to[R,l_={$\SI{6}{\ohm}$}] (7.0,0) coordinate(G);
  \draw (C) to[R,l_={$\SI{2}{\ohm}$}] (F);
  \draw (D)--(G);
  \foreach \x/\y/\n in {1.55/4.15/1,5.25/4.15/2,1.55/1.25/3,5.25/1.25/4}{%
    \draw[->,loopcol,line width=1.05pt] (\x,\y) arc(0:305:.34);
    \node[loopcol,fill=white,inner sep=1pt] at(\x,\y){\footnotesize$I_\n$};}
\end{circuitikz}}
